\chapter{Project Context and State of the Art}
\label{chap:context}
\begin{spacing}{1.2}
\minitoc
\thispagestyle{MyStyle}
\end{spacing}
\newpage

\section{Introduction}

This chapter first presents the host organisation in which this project
was carried out, then reviews the state of the art in EU CBAM regulation
and in agentic, tool-using language models, critiques that state of the
art against the specific demands of high-stakes numeric compliance work,
and formulates the resulting problem. It then reviews the existing
methodologies used for compliance-document automation and their limits,
and closes by presenting this project's proposed solution objectives,
research questions, and scope.

% ─────────────────────────────────────────────
\section{Host Organisation}

This project was carried out at \textbf{ninety. Startup House}, an
accelerator and incubator based in Sfax, Tunisia (Rte Mahdia KM~1.5, El
Bosten).

\begin{figure}[H]
\centering
\includegraphics[width=7.5cm]{ninety-logo.png}
\caption{ninety. Startup House logo.}
\label{fig:ninety_logo}
\end{figure}
\FloatBarrier

\subsection*{Mission and Vision}

ninety.'s stated mission is to assist innovative SMEs and startups
throughout their entrepreneurial journey, by offering a dynamic working
environment, personalised support, and the resources needed to turn an
idea into a functioning business. Its longer-term vision is to become a
primary innovation catalyst for the Sfax region, building an ecosystem in
which startups can not only survive but grow, and in doing so inspire
further entrepreneurship locally.

\subsection*{Core Values}

The organisation identifies five values as central to how it operates:

\begin{itemize}[label=\textbullet]
    \item \textbf{Innovation} fostering creativity and originality in the
    startups it hosts.
    \item \textbf{Collaboration} encouraging teamwork between startups,
    mentors, and investors.
    \item \textbf{Community} building a supportive network oriented
    toward shared success.
    \item \textbf{Excellence} pursuing continuous improvement and
    quality.
    \item \textbf{Ethics} emphasising transparency, accountability, and
    mutual respect.
\end{itemize}

\subsection*{Services}

Alongside its acceleration and incubation programmes, ninety. offers
coworking desks, private offices, dedicated training and study spaces,
meeting rooms for presentations and workshops, and business
domiciliation (virtual office) services for early-stage companies that
do not yet need a physical office of their own.

It is within this startup-support environment surrounding company
building and entrepreneurship that the present project's own subject
matter --- automating a specific, high-stakes regulatory compliance task
with agentic AI --- was defined and carried out.

% ─────────────────────────────────────────────
\section{State of the Art}

\subsection{EU CBAM as a Moving Regulatory Target}

Regulation (EU) 2023/956 establishes the Carbon Border Adjustment
Mechanism and its Annex~I/Annex~II scope~\cite{eu2023956}; Directive
2003/87/EC, Art.~10a(1a), defines the phase-out schedule for the
free-allocation benchmark deduction (2.5\% phased out in 2026, rising to
100\% by 2034)~\cite{eudirective200387ec}; Commission Implementing
Regulations 2025/2621 and 2026/1412 publish the yearly country/route
default emission values, markup schedule, and EU ETS free-allocation
benchmarks used in this project's reference data~\cite{eu20252621,eu20261412}.
The Joint Research Centre's JRC134682 report~\cite{jrc134682} is the
underlying scientific basis for these values. CBAM entered its
\textbf{definitive phase on 1~January 2026}: declarations became
\textbf{annual} rather than quarterly (the first annual declaration,
covering 2026 imports, is due 30~September~2027), and a binding
\textbf{50-tonne-per-year exemption} applies below which no certificate
obligation arises~\cite{eu2023956}. What distinguishes this regulatory
domain from a typical integration target is that the ``constants'' a
naive implementation would reach for are neither constant nor available
in one place: default emission values are published per exporting
country, carry a yearly markup that itself changes, and are periodically
revised by implementing acts~\cite{eu20252621} --- a requirement that
shaped the data model (Section~\ref{sec:calc_engine_cbam}) more than any
purely technical consideration did. A system that hardcodes these figures
in application code is wrong the moment a new implementing regulation is
published; the data must be modelled as a \textbf{versioned reference
table}, not a constant.

\subsection{Tool-Using LLMs and the ReAct Paradigm}
\label{sec:react_paradigm_cbam}

LLMs capable of multi-step reasoning and external tool use have enabled
agents that autonomously decompose complex goals into executable
steps~\cite{russell2020ai,wang2023agentsurvey}. Toolformer showed that a
model can learn when and how to call external APIs~\cite{schick2023toolformer},
and Gorilla that API-call accuracy depends heavily on how the tool
interface is presented to the model~\cite{patil2023gorilla}; AgentBench
found that open-weight models trail commercial ones most sharply on
exactly this kind of multi-turn, instruction-following
work~\cite{liu2024agentbench} --- a gap directly relevant to a project
restricted to locally hosted models. The \textbf{ReAct} framework~\cite{yao2022react}
grounds LLM reasoning in external reality by interleaving reasoning traces
with action steps --- calls to external tools or APIs. At each step the
agent produces a \textit{thought}, selects an \textit{action} (a tool call
with structured arguments), and receives an \textit{observation} (the
tool's result). This cycle repeats until sufficient information has been
gathered, or the agent escalates. Figure~\ref{fig:react_loop_cbam}
illustrates the loop as implemented in this project's orchestrator.

\begin{figure}[H]
\centering
\begin{tikzpicture}[
  node distance=1.2cm and 2.6cm,
  thought/.style={rectangle, rounded corners=4pt, draw=blue!60, fill=blue!8,
                  text width=3.4cm, align=center, font=\small, minimum height=0.8cm},
  action/.style={rectangle, rounded corners=4pt, draw=orange!70, fill=orange!8,
                 text width=3.4cm, align=center, font=\small, minimum height=0.8cm},
  obs/.style={rectangle, rounded corners=4pt, draw=teal!60, fill=teal!7,
              text width=3.4cm, align=center, font=\small, minimum height=0.8cm},
  env/.style={rectangle, rounded corners=4pt, draw=gray!60, fill=gray!8,
              text width=3.6cm, align=center, font=\small, minimum height=2.4cm,
              dashed},
  io/.style={rectangle, rounded corners=4pt, draw=green!60!black, fill=green!7,
             text width=3.1cm, align=center, font=\small, minimum height=0.8cm},
  arrow/.style={-{Stealth[length=5pt]}, thick},
]
\node[io] (task) {Declaration line to process};
\node[thought, below=0.9cm of task]    (thought) {\textbf{Thought}\\What is known so far / what is missing};
\node[action,  below=1.1cm of thought] (action)  {\textbf{Action}\\One of 7 tool calls, model-chosen};
\node[obs,     below=1.1cm of action]  (obs)     {\textbf{Observation}\\Tool result fed back};
\node[io, text width=4.6cm, font=\small\selectfont, below=0.9cm of obs] (answer)
  {\footnotesize\texttt{write\_declaration\_line}\\ or \\ \footnotesize\texttt{flag\_for\_human\_review}};
\node[env, right=2.8cm of action] (env)
  {\textbf{Odoo (read-only) +}\\\textbf{Calculation Engine}\\[5pt]
   \footnotesize extract\_field\\ lookup\_default\_value\\
   check\_consistency\\ calculate\_emissions, \ldots};
\draw[arrow] (task)    -- (thought);
\draw[arrow] (thought) -- (action);
\draw[arrow] (action)  -- (obs);
\draw[arrow] (obs)     -- (answer);
\draw[arrow, bend left=50]
  (obs.west) to node[left, font=\scriptsize, text width=1.3cm, align=center]
  {next\\reasoning\\step} (thought.west);
\draw[arrow] (action.east) -- node[above, font=\scriptsize] {tool call} (env.west);
\draw[arrow] (env.south)   -- ++(0,-0.5) -| node[pos=0.75, above, font=\scriptsize] {result}
  (obs.east);
\node[font=\scriptsize, text=gray, right=0.1cm of answer] {(single validated write)};
\node[font=\scriptsize, text=gray, right=0.1cm of task]   {(job start)};
\node[draw=blue!30, dashed, rounded corners=6pt,
      fit=(thought)(action)(obs), inner sep=8pt,
      label={[font=\footnotesize\itshape, blue!60]above:Agent Loop (max.\ 8--10 iterations)}] {};
\end{tikzpicture}
\vspace{0.3cm}
\caption{The ReAct-style agent loop~\cite{yao2022react} as implemented for
CBAM declaration processing.}
\label{fig:react_loop_cbam}
\end{figure}
\FloatBarrier

The distinguishing property of this execution model is that the loop
itself carries no knowledge of which tool ``should'' come next --- that
decision is made by the model at every turn, from the current conversation
state, not by a step counter or a branch on a previous tool's name. A
fixed pipeline that always calls ``extract quantity, then emissions, then
calculate'' will, on a document where the emissions figure is genuinely
absent, either fail outright or silently proceed with a placeholder; an
agent that decides its own order can instead request missing data, retry a
different extraction strategy, or escalate.

\paragraph{What ``ReAct-style'' means precisely here.} This project
borrows ReAct's interleaved reasoning-and-acting \textit{pattern}, not
its original \textit{mechanism}. Table~\ref{tab:react_vs_project} sets out
the difference.

\begin{table}[H]
\centering
\setstretch{1.05}
\small
\renewcommand{\arraystretch}{1.1}
\begin{tabularx}{\textwidth}{|L{3.0cm}|X|X|}
\hline
\textbf{Aspect} & \textbf{Original ReAct~\cite{yao2022react}} & \textbf{This project} \\
\hline
Action mechanism & Harness parses literal \texttt{Thought:}/\texttt{Action:} text markers & Native function calling: the model returns a structured tool call \\
\hline
Reasoning step & Part of the parsed output format & Prompted, never enforced or parsed \\
\hline
Turn skipping reasoning & Malformed output & Accepted unchanged \\
\hline
Role of the reasoning trace & Drives control flow & Counted for observability only \\
\hline
\end{tabularx}
\caption{Original ReAct versus this project's ReAct-style loop.}
\label{tab:react_vs_project}
\end{table}
\FloatBarrier

Not enforcing the reasoning step is deliberate: checking a turn's
\textit{content} is one step away from the code deciding what a turn
should contain, which would reintroduce the orchestrator-side ordering
logic this design rules out. Instead, the loop counts every tool-calling
turn and, separately, those whose content stated reasoning first, and
reports both through the job status endpoint. That lets the report say,
from evidence rather than assumption, how often the model actually
behaves in the ReAct pattern.

\subsection{Hallucination Risk in High-Stakes Numeric Extraction}

Hallucination surveys~\cite{ji2023hallucination,huang2023hallucination} document the
general tendency of generative language models to produce fluent,
plausible content unsupported by their input --- a risk that becomes a
compliance-integrity problem the moment the fabricated content is a number
that determines a certificate liability. Mallen et al.~\cite{mallen2023trust}
additionally show that LLM recall degrades specifically for rare, sparsely
represented facts relative to popular ones --- directly analogous, in this
project's domain, to a regulatory default value for an
infrequently-traded country/route combination being exactly the kind of
fact a general-purpose model is least reliable at recalling from
parametric memory alone. This motivates routing every default value
through a deterministic database lookup rather than ever trusting the
model to recall or compute one itself.

% ─────────────────────────────────────────────
\section{Critique of the State of the Art}

Taken individually, none of the three bodies of work reviewed above was
designed for the exact situation this project sits in, and each leaves a
gap that matters once the output is a compliance number rather than a
chatbot answer.

\textbf{The ReAct paradigm assumes, but cannot verify, that reasoning
actually happens.} Its central claim is that interleaving explicit
reasoning with action improves grounding and reduces error-propagation,
but that claim rests entirely on the model choosing to produce the
reasoning in the first place. Nothing in the paradigm itself, and nothing
in the native function-calling interfaces that have since replaced its
original text-parsing mechanism, gives an orchestrator a reliable way to
confirm that a ``thought'' actually preceded an ``action'' rather than
being an after-the-fact rationalisation. The only two ways to check would
be to constrain the model's output format, which defeats the point of
letting it call tools freely, or to simply trust it. In short, the
paradigm is a prompting discipline, not a verifiable guarantee.

\textbf{General hallucination-mitigation techniques target the wrong
granularity for this problem.} Retrieval
augmentation~\cite{lewis2020rag}, self-consistency voting across multiple
samples~\cite{wang2023selfconsistency}, and learned fact-checking
classifiers --- the mitigations most commonly proposed against
hallucination~\cite{huang2023hallucination} --- are all built for open-domain factual question
answering, where the model has to recall or retrieve a fact from a
large, unbounded knowledge base. This project's extraction problem is
narrower and more tractable than that: the value to be grounded is
already sitting, in full, inside one specific document the model has
already been shown in its own context window. The general literature
does not distinguish this ``verify against a document already in
context'' case from the harder open-domain case it was designed for,
which leaves the specific, cheaper check this project actually relies on
(Chapter~\ref{chap:implementation}) unaddressed by name.

\textbf{The rare-fact-recall finding motivates a fix the literature does
not itself supply.} Mallen et al.'s own proposed remedy is adaptive
retrieval, which assumes an underlying corpus large and general enough to
retrieve from. CBAM's regulatory default values do not fit that setting:
they form a small, precisely-scoped, versioned table with a handful of
dimensions (sector, route, country, validity period), not an open corpus.
A lookup this constrained does not need retrieval-augmented generation at
all --- a deterministic query answers it exactly, every time, with no
retrieval-relevance uncertainty --- but the literature reviewed above
never separates this bounded-lookup case from the general recall problem
it studies.

None of these three gaps rules out using a tool-using LLM for this task;
each just means the general-purpose techniques in the literature need
adapting, or replacing with a narrower, purpose-built mechanism, rather
than applied as published. What remains is to state precisely where an
agentic, tool-using LLM can be trusted, and where it cannot, for a task
whose output is a real financial liability.

% ─────────────────────────────────────────────
\section{Problem Statement}

This regulatory design has an asymmetric commercial consequence for
exporters outside the EU. Tunisia exports steel, aluminium, and phosphate
products to the European Union, and is directly affected by CBAM
reporting requirements that originate on the EU-importer side: a
Tunisian mill's own paperwork becomes an input to a foreign compliance
pipeline, without any change on the exporter's part. That is why the
extraction layer is built against realistic, heterogeneous document
formats rather than a single clean template --- the documents an
importer actually receives from a diverse supplier base are never
uniform.

\subsection*{The Compliance-Automation Dilemma}

A CBAM declaration line requires combining several pieces of information
that are rarely all present, in a consistent format, in a single
document: a quantity, a production route, an exporting country, and
either a measured emissions figure or a default value looked up from the
versioned reference table described above. Reading and reconciling this
information from real supplier paperwork is exactly the kind of
unstructured-text task large language models are good at. The problem is
what happens next --- the number that comes out the other end determines
a \textit{real certificate liability}. An LLM that hallucinates a
plausible-looking emissions figure, or silently drops a required field
under multi-argument tool-call construction (both observed in this
project's own evaluation, Section~\ref{sec:results_cbam}), produces a
wrong compliance number that looks exactly like a correct one. Any
system that lets a language model touch the arithmetic directly inherits
that risk unconditionally.

\subsection*{The Workflow-vs-Agent Gap}

Most existing document-automation tooling --- scripted extraction
pipelines with an LLM call bolted onto one step --- encodes a
predetermined execution path: extract field A, then field B, then
compute. This is brittle against the real variability of supplier
documents, and it conflates two concerns that should really be kept
separate: deciding what to do next given what has been read so far (a
reasoning problem, suited to an LLM), and computing a number from
validated inputs (a deterministic problem that must never depend on a
model's judgment). Left unresolved, this conflation exposes a business
to three compounding risks:

\begin{itemize}[label=\textbullet]
    \item \textbf{Silent numeric hallucination:} a fabricated emissions
    value is indistinguishable, in the output, from a genuinely extracted
    one, unless the system is explicitly designed to demand and verify
    provenance.

    \item \textbf{Brittleness to document variability:} a fixed step
    sequence has no recovery strategy when an expected field is genuinely
    absent from a given supplier's document, beyond failing the whole run.

    \item \textbf{Silent drift against regulatory data:} default emission
    values and benchmarks that are hardcoded, rather than versioned and
    cited, become quietly wrong the moment an implementing act revises
    them.
\end{itemize}

A system that wants to be robust to document variability, without also
inheriting an LLM's arithmetic unreliability, has to separate these
concerns architecturally rather than merely prompt around the risk. That
is the precise problem this project is designed to solve.

% ─────────────────────────────────────────────
\section{Existing Methodologies for This Problem and Their Limits}
\label{sec:related_work_cbam}

Three families of existing approaches apply to this problem.

\textbf{Robotic process automation (RPA).} RPA tools operate on the user
interface of other systems the way a human operator would, replaying a
recorded procedure~\cite{vanderaalst2018rpa}. They are fast and
predictable on the exact layout they were built against, but the RPA
literature itself lists brittleness to process and interface change among
the field's open challenges~\cite{syed2020rpa}. Applied to supplier
documents, a supplier reformatting a mill certificate can silently break
every downstream field mapping, and the robot cannot tell ``this field is
absent'' from ``this field moved.''

\textbf{Learned document-understanding models.} Layout-aware models such as
LayoutLM~\cite{xu2020layoutlm} and OCR-free models such as
Donut~\cite{kim2022donut} extract fields from forms and receipts far more
robustly than templates. They are, however, trained extractors: they
require labelled training data for each document type, produce a single
pass of predictions with no ability to ask for missing information, and
have no notion of a downstream calculation they must stay separate from.
No labelled CBAM corpus exists to train them on.

\textbf{Generative LLM extraction.} Generative LLMs can extract structured
fields from free text with little or no task-specific training, and this
is now an active research area in its own right~\cite{xu2023genie}. Most
such pipelines, however, apply one extraction call per document, with no
iterative reasoning and no deterministic boundary between reading a value
and computing with it; when a downstream figure is needed, it is often
computed inside the same model call that did the reading. Tool-use
benchmarks also show that open-weight models are markedly less reliable
at multi-turn, instruction-following agent work than commercial
ones~\cite{liu2024agentbench}, so published results obtained with large
hosted models do not transfer directly to the small, locally hosted
models a data-sovereign deployment would use. Finally, documents fed to
an LLM are untrusted input: indirect prompt injection through processed
content is a demonstrated attack on LLM-integrated
applications~\cite{greshake2023injection}, which matters when the output
is a regulatory figure.

\begin{table}[H]
\centering
\setstretch{1.05}
\begin{tabularx}{\textwidth}{|L{3.0cm}|X|X|}
\hline
\textbf{Aspect} & \textbf{Existing Approaches} & \textbf{This Project} \\
\hline
Execution model &
Fixed pipeline or single-shot LLM call &
Autonomous ReAct-style loop; model decides tool order/count at runtime \\
\hline
Compliance arithmetic &
Often computed inline by the LLM or rigid template logic &
Deterministic Python, zero LLM involvement, independently unit-tested \\
\hline
Regulatory data &
Frequently hardcoded or embedded in prompts &
Versioned reference table, never hardcoded, every value cited to a primary source \\
\hline
Hallucination defense &
Rarely addressed explicitly &
Multi-layer: plausibility bounds, untrusted-content framing, verbatim source-quote grounding \\
\hline
Model selection &
Assumed from vendor/model-card claims &
Empirically benchmarked, with a documented self-correction of an initial, biased ranking \\
\hline
\end{tabularx}
\caption{Existing compliance-automation approaches vs.\ this project's design.}
\label{tab:existing_vs_project}
\end{table}
\FloatBarrier

Table~\ref{tab:existing_vs_project} positions this project's own choices
against these limits directly. No widely available compliance-document
tool separates extraction (agentic) from calculation (deterministic) as a
first-class constraint; few publish a harness-matched benchmark of which
locally-hostable model is reliable enough before committing to one; and
CBAM's versioned, country/route-dependent default-value regime combined
with an agentic extraction loop has, to the author's knowledge, no
existing open implementation to compare against. This is the space the
proposed solution below occupies.

% ─────────────────────────────────────────────
\section{Proposed Solution}

This project resolves the gaps identified above architecturally, rather
than by prompting around them. A deterministic, LLM-free calculation
engine is kept categorically separate from a genuinely agentic tool-use
loop for document extraction, and candidate models are benchmarked
empirically against the project's own production validation logic ---
rather than trusted on the strength of a published card, or a generic
tool-calling benchmark whose scoring rule may not match how the real
system actually consumes a model's output. Concretely, the project aims
to design and implement a genuinely agentic tool-use system for CBAM
declaration automation, paired with a deterministic
compliance-calculation core. Three objectives guide its development:

\begin{itemize}[label=\textbullet]
    \item \textbf{O1.} Implement a deterministic, unit-tested calculation
    engine for CBAM embedded emissions and certificate liability, with zero
    involvement from any language model in producing the compliance number
    itself.

    \item \textbf{O2.} Design and implement a genuine agentic tool-use
    loop --- the model decides its own tool-call order and count, with no
    hardcoded step sequence in the orchestrating code --- for supplier
    document extraction and reconciliation.

    \item \textbf{O3.} Empirically benchmark candidate locally-hosted
    language models for tool-calling reliability and end-to-end
    field-extraction accuracy, rather than assuming capability from a
    model's published card.
\end{itemize}

Two research questions guide the investigation:

\begin{itemize}[label=\textbullet]
    \item \textbf{RQ1.} How reliably can a small, locally-hosted language
    model execute a self-directed, multi-step tool-use loop for structured
    regulatory data extraction, and does a raw tool-calling benchmark
    actually predict that reliability?

    \item \textbf{RQ2.} What architectural and prompting safeguards are
    needed to prevent a tool-using language model from corrupting a
    compliance-critical numeric output, given that document text fed to
    the model is untrusted external input?
\end{itemize}

The system covers \textbf{steel and aluminium imports only}
(Annex~II's direct-emissions-only sectors), uses \textbf{synthetic
supplier documents} (no real company data was available), and targets a
\textbf{two-machine deployment} meshed over a private VPN. Real-time
monitoring, sectors beyond steel/aluminium, and a production-grade
multi-tenant deployment fall outside the scope of this project. As
Table~\ref{tab:existing_vs_project} above summarises, each of these
objectives answers a specific limitation of the existing approaches
already reviewed, rather than a generic ambition to ``add AI'' to
compliance work.

% ─────────────────────────────────────────────
\section{Conclusion}

This chapter has situated the project at ninety. Startup House, reviewed
CBAM's shift toward annual, versioned-data declarations together with the
agentic-AI and hallucination literature it depends on, and shown where
that literature falls short of the specific demands of a high-stakes
numeric compliance task. It then formulated the resulting
compliance-automation dilemma, reviewed the RPA-based and single-shot
LLM approaches currently used against it and their shared limits, and
presented this project's own objectives, research questions, and scope as
a direct response to those limits. These constraints --- a deterministic
core, a genuinely autonomous extraction loop, and an empirical model
choice --- directly shape the design presented in the following chapter.
